\chapter{Literature Review}\label{Chap:1}

\section*{Introduction}

\section{Malaria Drug Discovery Background}
Malaria continues to represent a substantial global public health concern, exerting a profound impact on regions where it is endemic. The intricate biology of the malaria parasite, involving its lifecycle and interactions with Anopheles mosquitoes, underscores the necessity for a comprehensive understanding of the dynamics of the disease. Recent research underscores the influence of climatic factors on malaria distribution, particularly in the highlands of East Africa, where defying traditional boundaries signifies a potential threat to previously unaffected populations \cite{Chaves2010}. Advancements in computational biology have yielded novel insights into the intricacies of malaria transmission, facilitating targeted interventions and policy recommendations \cite{Chasapi2019}. A foundational expertise in malaria biology is essential not only for academic purposes but also for devising effective strategies to combat the disease, thereby enhancing public health outcomes in vulnerable regions.

\subsection{Overview of malaria biology}
The persistent impact of malaria on global health is underscored by its historical prevalence and ongoing threat, particularly in tropical and subtropical regions. Malaria, a potentially lethal disease caused primarily by the Plasmodium falciparum parasite, imposes a substantial burden on public health systems, resulting in elevated mortality rates, particularly among vulnerable populations such as children and pregnant women. In response, concerted efforts have been made to combat the disease, propelling advancements in genomic research aimed at identifying novel biological targets and drug candidates, which are essential for the development of effective treatment and prevention strategies \cite{Birkholtz2006}. Furthermore, shifts in climatic conditions have facilitated the resurgence of malaria in regions previously considered low-risk, such as the highlands of East Africa, thereby sparking debates regarding the role of climate change in its geographic expansion \cite{Chaves2010}. This multifaceted challenge underscores the necessity for a comprehensive understanding of the biological and environmental factors that influence malaria's persistence as a global health concern.

\subsubsection{Different species of Plasmodium}
Malaria is a parasitic disease instigated by protozoans of the \emph{Plasmodium} genus, which are members of the \emph{Apicomplexa} family. Transmission to humans occurs through the bite of an infected female Anopheles mosquito. Historically, four Plasmodium species were acknowledged as human pathogens. However, recent scientific investigations have uncovered a fifth species that had previously gone undetected \cite{CoxSingh2008,Singh2004}, thereby adding complexity to the epidemiological landscape of malaria. This finding signifies an increase in the number of Plasmodium species known to infect humans, now totaling five (see \Cref{fig:mal_type}).
\begin{figure}[ht]
	\includegraphics[scale=0.92]{malaria_type.pdf}
	\caption[Different species of plasmodium]{Different species of plasmodium : (i) Plasmodium falciparum, (ii) Plasmodium vivax, (iii) Plasmodium oval, (iv) Plasmodium malariae, (v) Plasmodium knowlesi}
	\label{fig:mal_type}
\end{figure}
Thus we have :
\tcbset{fonttitle=\bfseries\upshape,
	fontupper=\slshape, colback=green!8!white,colframe=green!55!black, separator sign dash, breakable, width=1.04\linewidth}
\begin{tcolorbox}
	\begin{itemize}
		\item \textbf{\textcolor{codered!80}{Plasmodium Falciparum}} (\emph{Pf}) : This species is responsible for the most severe form of the disease, which frequently results in fatal outcomes. The species under consideration is present in the tropical zones of the African, Latin American, and Asian continents. This species has been observed to develop resistant forms, which have been documented in the context of existing antimalarial treatments.
		\item \textbf{\textcolor{codeblue!70}{Plasmodium Vivax}} (\emph{Pv}) : The phenomenon under consideration has been observed to co-exist with \emph{Pf} in a variety of geographical regions worldwide, with a notable prevalence in South America and Asia. It has been identified as a causative agent of mild, benign fever. The prevalence of cases involving the consumption of \emph{Pv}, which have been reported to result in severe, even fatal, outcomes, underscores the potential risks associated with this species \cite{Getachew2015,Barnadas2015,Marfurt2008}. 
		
		\item \textbf{\textcolor{codebrown!60}{Plasmodium Oval}} (\emph{Po}) : The presence of the specie has been documented in the intertropical region of Africa, as well as in specific regions within the Pacific, with a notable prevalence in areas where the presence of the \emph{Pv} is absent. \emph{Po} does not result in immediate mortality; however, it has been observed to induce a state of reviviscence that persists for a duration ranging from four to five years.
		
		\item \textbf{\textcolor{codegreen!65}{Plasmodium Malariae}} (\emph{Pm}) :  The species exhibits a scattered geographical distribution. Although \emph{Pm} is not lethal, it has the potential to cause relapses twenty years after the initial infection.
		
		\item \textbf{\textcolor{codepurple!70}{Plasmodium Knowlesi}} (\emph{Pk}) : It is evident that infections in humans due to Pk are a recent phenomenon. The distribution of this species is limited to the Asian continent. This is explained because this species has long been confused with \emph{Pm} \cite{CoxSingh2008,Singh2004}.
	\end{itemize}
\end{tcolorbox}	

\subsubsection{Life cycle of malaria parasite}
The complex interplay of malaria transmission is predominantly characterized by the interaction between the Anopheles mosquito and the human host. The various stages of parasite (see \Cref{fig:life_cycle}) development play a pivotal role in determining the organism's pathogenic potential. Upon the bite of an infected mosquito, the parasite known as a sporocyst is injected into the victim. This sporocyst rapidly migrates to the liver, where it initiates asexual reproduction. Ultimately, the merozoites disperse into the bloodstream. The malaria parasite displays a multifaceted life cycle, which is initiated through the interaction with an Anopheles mosquito and subsequently continues within the body of a vertebrate host. Following entry into the bloodstream, these merozoites invade red blood cells, where they undergo further multiplication, leading to the symptomatic manifestations of malaria. Recent studies have emphasized the post-transcriptional regulation of gene expression in the parasite, with a focus on the over \num{1000}-fold range. The presence of RNA-binding proteins has been identified, thereby suggesting the presence of a sophisticated mechanism of gene regulation that plays a crucial role in orchestrating the life cycle and survival of the parasite \cite{Bunnik2016}. It is imperative to comprehend the intricacies of this phenomenon to formulate targeted interventions against malaria \cite{Linzke2019}.
\begin{figure}[H]
	\includegraphics[scale=0.115]{Life_Cycle_of_the_Malaria_Parasite.pdf}
	\caption[Life cycle of the malaria parasite]{Life cycle of the malaria parasite : In the human host, the parasite Plasmodium introduces itself into the bloodstream during a mosquito blood meal (\textbf{\tcr{Step 1}}) and rapidly targets hepatocytes (\textbf{\tcr{Step 2}}), where it undergoes an exoerythrocytic schizogony over approximately 5–6 days. In the context of a hepatocyte infected with a parasite, a single sporozoite will undergo differentiation into a schizont, yielding approximately \num{40000} merozoites. These merozoites will then be released into the circulation. These merozoites invade erythrocytes (\textbf{\tcr{Step 3}}), thereby initiating the 48-hour asexual erythrocytic cycle. Each trophozoite matures into a schizont that liberates $4–30$ new merozoites upon erythrocyte rupture. A subset of merozoites undergoes differentiation into sexual forms (\textbf{\tcr{Step 4}}), thereby giving rise to microgametocytes and macrogametocytes within circulating red cells. Subsequently, during a blood meal (\textbf{\tcr{Step 5}}), an Anopheles mosquito ingests these gametocytes. Within the midgut lumen, the gametocytes undergo exflagellation, fusion, and fertilization to form a diploid zygote (\textbf{\tcr{Step 6}}). The zygote subsequently transforms into a motileookinete, which traverses the midgut epithelium and establishes an oocyst on its basal surface (\textbf{\tcr{Step 7}}). It has been determined that during the 8–15 days of sporogony, each oocyst generates approximately \num{3000} sporozoites. These sporozoites subsequently migrate to the mosquito's salivary glands (\textbf{\tcr{Step 8}}), thereby completing the cycle and rendering the vector capable of transmitting infection to a new human host.}
	\label{fig:life_cycle}
\end{figure}

\subsubsection{Pathophysiology of malaria}
The profound pathophysiological effects of malaria serve to underscore the disease's considerable impact on human health. As the disease progresses, particularly in cases of \emph{Pf} infection, systemic inflammation and the sequestration of parasitized erythrocytes become pivotal. It has been established that these processes are responsible for severe manifestations, such as cerebral malaria, which is characterized by altered neurological functioning. Notably, studies indicate that these substances contribute to the onset of coma in the brain. In the placenta, they have been shown to contribute to low birthweight and preterm labor, and to increase the risk of abortion and stillbirth. In the context of \emph{Pv}, recent advancements in genetics have facilitated a more nuanced comprehension of its mechanisms. Nevertheless, challenges persist due to the constrained capacity of laboratory studies for this species \cite{Auburn2016}. As demonstrated in the research by \cite{Getachew2015}, both the inflammatory response and the genomics of pathogens play crucial roles in determining disease severity and the interindividual susceptibility observed among affected populations.
 
\subsubsection{Mechanisms of disease and immune response in infected individuals}
The intricate interplay between malaria pathogens and the host immune response unveils significant insights into disease mechanisms.  As Plasmodium species invade erythrocytes, they employ antigenic variation to evade immune detection. This evasion is notably achieved through the expression of variant surface proteins, such as PfEMP1, which complicates the host's ability to mount an effective immune response. This dynamic is further influenced by the splenic components that are essential for the processing of variant antigens. This phenomenon has been demonstrated in studies of \emph{Pk} infections, where specific antibodies facilitate variant antigen switch events in vivo. As the immune system interacts with these constantly changing antigens, antibodies evolve, thereby enhancing the host's defensive capabilities over time. However, the critical challenge of the essential delay in antibody production underscores the necessity of immunization strategies in combatting malaria. Without prior exposure, the body struggles to respond swiftly to novel infections \cite{GALINSKI2017}. 

\subsection{Pharmaceutical drug discovery process}
In the medical, biotechnology, and pharmacology fields, the process of drug discovery entails the identification of novel potential medications. The contemporary process of drug discovery involves the identification of screening hits, the application of medicinal chemistry, and the refinement of these hits to enhance their affinity, selectivity (to minimize side effects), efficacy/potency, metabolic stability (to extend the half-life), and oral bioavailability. Once a compound that meets all these criteria has been identified, it progresses to drug development before entering clinical trials. This contemporary approach to drug discovery is ordinarily a capital-intensive endeavor, necessitating substantial investments from pharmaceutical companies and national governments. The culmination of the drug discovery process is the acquisition of a patent for the drug in question. The pharmaceutical company must conduct costly Phase I, II, and III clinical trials, and many drugs fail to succeed in these trials. It is imperative to acknowledge the pivotal role that small companies play in this context. Frequently, these enterprises find themselves in a position where they must relinquish their rights to larger firms that are better equipped to undertake the necessary trials. The identification of drugs that achieve commercial or public health success is characterized by a multifaceted interaction among investors, industry, academia, patent laws, regulatory exclusivity, marketing, and the necessity of balancing confidentiality with communication \cite{Helleboid2014}.

\begin{figure}[H]
	\includegraphics[scale=0.3]{dd_process}
	\caption{Drug discovery process}
	\label{fig:dd_process}
\end{figure}
 The main steps of the drug discovery process (see \Cref{fig:dd_process}) are listed below:
 \begin{itemize}
 	\item \textbf{Target Validation} : The process of target validation in the pharmaceutical industry entails the evaluation of cellular or molecular structures associated with diseases. The term \emph{target} remains a subject of debate; however, the industry distinguishes between established and new targets. The established targets have a substantial scientific literature documenting their role in physiology and disease, requiring less investment for therapeutic development and offering insights on chemical feasibility. New targets lack this established background but are currently pursued in discovery efforts, often representing newly identified proteins or those with recently elucidated functions. As of 2011, \num{435} human genome products had been identified as therapeutic drug targets for FDA-approved medications \cite{Elion1993}.
 	
 	\item \textbf{Target to lead} : The process of drug discovery is commonly initiated with high-throughput screening (HTS) of chemical libraries, with the objective of identifying compounds that exert an effect on the target. Substances that successfully complete the preliminary screening process are then subjected to cross-screening, a methodical evaluation that aims to substantiate their selectivity and mitigate potential risks of off-target toxicity. The elimination of \emph{pan-assay interference compounds} is initiated at this stage. In instances where a plurality of active compounds exhibit shared chemical characteristics, practitioners of medicinal chemistry may develop pharmacophores \cite{Black2025}. Subsequently, they enhance lead compounds through iterative structure-activity relationship (SAR) studies to improve potency, selectivity, and ADMET properties.
 	
 	\item \textbf{Clinical trials} : The design of new pharmaceuticals clinical trials is generally divided into four phases, although individual trials may span multiple phases, such as combined phase I/II or phase II/III trials. Consequently, a more systematic approach might involve the classification of these studies as either early phase or late phase \cite{FDA2025}. The process of developing a new pharmaceutical compound typically traverses all four phases over the course of several years. The successful navigation of Phases I, II, and III by a pharmaceutical agent typically results in its approval for public use by the national regulatory authority. Phase IV encompasses studies conducted subsequent to approval.
 	\begin{enumerate}
 		\item \emph{Phase I} trials were formerly referred to as "first-in-man studies," but the field generally transitioned to the gender-neutral language phrase "first-in-humans" in the 1990s. These trials represent the initial stage of testing in human subjects. These studies are conducted with the objective of ascertaining the maximum tolerable dosage of the pharmaceutical agent in question. The primary endpoint of these investigations is to identify the highest dosage that can be administered to a subject without incurring adverse effects that would compromise the patient's well-being or safety. Typically, a cohort of 2–100 healthy volunteers is recruited for this purpose. 
 		
 		\item \emph{Phase II} : Subsequent to determining a dose or range of doses, the subsequent objective is to evaluate whether the pharmaceutical agent exhibits any biological activity or effect. These trials are conducted on larger groups (100–300) and are designed to assess the efficacy of the drug, as well as to continue Phase I safety assessments in a larger group of volunteers and patients.
 		
 		\item \emph{Phase III} : The objective of this phase is to assess the effectiveness of the new intervention and, consequently, its value in clinical practice. Phase III studies are randomized controlled multicenter trials on large patient groups (\num{300}–\num{3000} or more, depending on the disease or medical condition studied) and are aimed at providing a definitive assessment of the efficacy of the drug in comparison with the current \textbf{gold standard} treatment.
 		
 		\item \emph{Phase IV} is also referred to as a postmarketing surveillance trial or, more colloquially, a confirmatory trial. These trials entail the safety surveillance (pharmacovigilance) and ongoing technical support of a drug following its authorization for commercialization (e.g., subsequent to approval under the FDA Accelerated Approval Program) \cite{Lancet2009}.
 	\end{enumerate}
 \end{itemize}
 
\subsection{Current treatments}
Notwithstanding substantial progress in the domain of medical science, infectious diseases continue to present considerable challenges to global health, particularly in tropical and subtropical regions, where access to healthcare resources exhibits significant variability. The repercussions of malaria extend beyond the realm of public health, encompassing socioeconomic development and impeding the capacity of healthcare systems to function effectively. The aforementioned complexity is further compounded by cultural factors that influence treatment preferences. This phenomenon is evident in regions where traditional medicine is often preferred due to its perceived efficacy or limited access to conventional biomedical treatments \cite{Gyasi2016}. Moreover, the emergence of drug resistance against conventional antimalarial therapies is indicative of broader concerns in the management of infectious diseases, akin to the challenges observed with antibiotics, where stagnation in new drug discovery and increasing resistance jeopardize treatment efficacy \cite{Barmpouni2022}. 

\subsubsection{Traditional treatments for malaria}
Historical methodologies employed in the treatment of malaria have placed significant emphasis on traditional remedies derived from plant sources \cite{Newman2007}, as well as on indigenous knowledge, predating the emergence of modern pharmacology by considerable periods of time. Quinine, a natural substance derived from the bark of the cinchona tree, was among the earliest and most significant treatments for malaria. It laid the foundation for the development of modern antimalarial drugs. Despite the advancements in the field of synthetic drug development, many communities persist in their utilization of herbal concoctions, driven by factors such as accessibility and entrenched cultural beliefs. However, the efficacy and safety of such traditional treatments vary widely and often lack rigorous scientific validation. While certain traditional practices may offer potential therapeutic benefits, the emergence of drug-resistant malaria strains highlights the urgent need for treatments that are subjected to rigorous scientific evaluation. Complementary and alternative medical practices, despite their prevalence, often exhibit outcomes that are no more efficacious than placebo when subjected to rigorous study methods. This observation underscores the significant challenges associated with the integration of these remedies into the realm of evidence-based medicine. Consequently, the reliance on traditional treatments must be balanced with advancements in clinically proven therapies to effectively combat malaria.

\subsubsection{Current advances in malaria treatment}
Recent advancements in malaria treatment have expanded beyond conventional indoor interventions to address the challenge posed by outdoor-biting mosquitoes, which maintain transmission even in areas with widespread use of bed nets and indoor spraying. Innovative tools such as the \emph{Mosquito Landing Box}, baited with human odors and treated with mosquitocidal agents, have demonstrated efficacy in attracting and killing outdoor host-seeking vectors, particularly Anopheles species responsible for malaria transmission. As demonstrated by field studies, the deployment of these devices has been shown to enhance the capture of mosquitoes when baited with natural human odors. Furthermore, the optimization of contact with effective insecticides, such as pirimiphos-methyl, has been observed to result in an increase in vector mortality. However, challenges persist in ensuring the rapid and sustained insecticidal action, given the brief visits of mosquitoes \cite{Matowo2013}. Moreover, the global fight against malaria is intricately linked to a broader array of global concerns regarding antimicrobial resistance. This underscores the necessity for a coordinated development and distribution of new diagnostics, vaccines, and treatments that are adapted to local conditions. As Bhagwandin asserts \cite{Bhagwandin}, effective governance and innovative incentive mechanisms are crucial to sustaining research and overcoming resistance hurdles.

\subsubsection{Antimalarial medecines}
In recent decades, there has been a shift in focus towards artemisinin-based combination therapies (ACTs), which have demonstrated efficacy against strains of Plasmodium falciparum that are resistant to other treatments. However, the emergence of resistance to these newer treatments underscores the urgent need for innovative solutions and ongoing research in the field of anti-malarial drug development \cite{Na‐Bangchang2009,Tse2019}.  The exploration of new candidates is underway, with a particular focus on the identification of compounds that target novel mechanisms of action to combat resistance and enhance treatment outcomes. This persistent endeavor is of paramount importance as malaria continues to represent a substantial public health concern on a global scale.  The global burden of malaria remains substantial, with millions of cases reported annually, particularly in sub-Saharan Africa, where the disease is most prevalent \cite{Kondrachine1997}. Addressing this challenge necessitates not only the development of new drugs but also a comprehensive approach that includes vector control, effective surveillance, and community engagement in prevention efforts.  The integration of these strategies is imperative to enhance the effectiveness of anti-malarial interventions and ensure sustainable disease control. Moreover, the integration of traditional medicinal knowledge may offer novel insights into the potential identification of natural products that could serve as novel anti-malarial agents \cite{Dwivedi2022}.


 
\subsection{Emerging resistance issues}
The efficacy of antimalarial medications is increasingly undermined by the complex strategies employed by malaria parasites to evade drug treatments. Genetic alterations in Plasmodium species have the potential to modify the targets of pharmaceutical drugs or enhance the parasites' capacity to expel therapeutic substances, thereby reducing the drugs' effectiveness. For instance, the resistance exhibited by parasites to artemisinin \cite{Folake2018,Kakde2023}, the prevailing treatment, has been linked to mutations in the Kelch13 gene. These mutations have been shown to influence the parasite's protein balance and its ability to recover from the adverse effects of the drug. Furthermore, the development of resistance to other antimalarials, including quinine derivatives and antifolates, is attributable to alterations in parasite transport proteins and enzymes that are essential for metabolic processes. This complicates the formulation of effective treatment plans. Recently, sophisticated techniques, including photoaffinity labeling, have played a pivotal role in elucidating these molecular changes by ascertaining the precise biological targets and resistance pathways of antimalarial drugs in live parasites. It is imperative to comprehend these mechanisms to develop novel therapeutic approaches that can counteract resistance and maintain malaria control efforts. The increasing resistance of malaria parasites and vectors to current interventions poses a significant challenge to existing treatment and control strategies \cite{Sinha2014}. Artemisinin-based combination therapies (ACTs), which were once the foundation of effective treatment, are now threatened by the rise of artemisinin resistance, especially in Southeast Asia. This development risks a reversal of progress made in recent decades. This situation necessitates the implementation of immediate modifications to therapeutic protocols, encompassing the exploration of triple artemisinin combination therapies (TACTs) and the evaluation of alternative drug candidates to ensure the maintenance of efficacy. Concurrently, the emergence of insecticide resistance in vector populations has been observed, thereby diminishing the efficacy of established interventions such as insecticide-treated nets and indoor residual spraying. Innovations such as gene drive technologies and next-generation insecticides hold great promise; however, they must be subjected to a thorough evaluation process to ensure their safety and long-term sustainability. The effectiveness of malaria control strategies is contingent upon the continuous implementation of genomic surveillance, a process that facilitates the expeditious identification of resistance patterns and the subsequent development of adaptive interventions. In the absence of integrated, dynamic strategies, attaining the WHO's ambitious 2030 reduction goals remains improbable. Moreover, recent advancements in genome-based technologies are facilitating the development of innovative approaches to address drug resistance, enabling the identification of new therapeutic targets and the creation of more effective treatment regimens. This collaborative endeavor among various stakeholders is imperative to ensure a comprehensive strategy to combat malaria and achieve long-term control and eradication goals \cite{Wells2010}.
\begin{figure}[H]
	\centering
	\includegraphics[scale=0.61]{antimalaria_resistant.pdf}
	\caption[Different parameters that contribute to antimalarial resistance]{Different parameters that contribute to antimalarial resistance \cite{Sinha2014}.}
	\label{fig:antimalaria_resistant}
\end{figure}

\section{Natural Products in Drug Discovery}
Natural products (NPs) are substances derived from natural sources such as plants, animals or minerals. They have been used in traditional medicine for centuries and have gained popularity in recent years as people seek natural alternatives to conventional medicines and other synthetic products \cite{NtieKang2019}. There are many different types of NPs, including herbal remedies, essential oils, dietary supplements and organic foods. Some NPs have been scientifically proven to have health benefits, while others have not been thoroughly studied or may be harmful.

\subsection{Natural products as drug}
Many pharmaceutical drugs are derived from NPs or have been inspired by natural compounds. For example, aspirin is derived from salicylic acid, which is found in the bark of willow trees, and the antimalarial drug artemether is derived from artemisinin, which is extracted from the plant \emph{Artimisia annua}, sweet wormwood, a herb used in traditional Chinese medicine. NPs can have a wide range of therapeutic benefits, including antimicrobial, anti-inflammatory, analgesic and immunomodulatory properties. They may also have fewer side effects than synthetic drugs, making them an attractive alternative for some people. However, developing NPs into drugs can be challenging due to issues such as low yields, difficulty in synthesising the compounds and variability in the composition of natural sources. The drug discovery process for NPs is given in \Cref{fig:np_drug_process}.
\begin{figure}[htbp]
	\centering
	\includegraphics[scale=0.57]{NP_process.pdf}
	\caption{Natural products drug discovery and development from plant}
	\label{fig:np_drug_process}
\end{figure}
This process fits into the traditional drug discovery process of target identification and validation, hit identification, hit to lead, lead to candidate and development. In addition, NPs can interact with other drugs and have potential side effects, so it's important to use them under the guidance of a healthcare professional. Overall, NPs have the potential to be valuable sources of new medicines, but rigorous scientific research is needed to fully understand their efficacy, safety and mechanisms of action.

\subsection{Successes of natural products and the untapped potential of African medicinal plants}
The exploration of the remarkable successes of natural products, especially within the captivating and fascinating context of African medicinal plants, reveals a vibrant and rich tapestry that is intricately woven with a wealth of traditional knowledge, immense biodiversity, and immense potential for groundbreaking, modern pharmacological applications. An extensive body of literature exists that highlights the considerable significance of these plants as vital and indispensable resources for healthcare, particularly in regions where access to conventional medicine is severely limited or even nonexistent. This underscores the critical importance of meticulously preserving, studying, and understanding these invaluable plants. They hold the extraordinary promise of unlocking new therapeutic options and significantly enhancing our understanding of medicinal chemistry and the intricate relationships between natural compounds and human health.

Obembe et \emph{al.} \cite{DikeI.2012}, underscores the effectiveness of Nigerian medicinal plants, which are often preferred over commercially available drugs. The chemical diversity inherent in Nigeria's flora presents a promising avenue for the discovery of novel lead compounds. However, the unsustainable harvesting practices pose a threat to the survival of certain species, underscoring the necessity for conservation efforts and the acknowledgement of traditional medicine's significance in healthcare. Ntie-Kang et \emph{al.} build upon this foundation by underscoring the substantial reservoir of phytochemicals present in African medicinal plants, particularly within the context of traditional medicine \cite{NtieKang2013a}. The study illustrates how local populations have historically relied on these plants for treating various medical conditions. The authors advocate for a collaborative approach involving traditional healers and scientists to validate the medicinal uses of these plants, thus facilitating the development of modern therapeutic agents. Additionally, Onguéné et \emph{al.} explore the anti-malarial properties of alkaloids and terpenoids derived from African medicinal plants \cite{AmoaOnguene2013}. The authors posit that the integration of traditional knowledge into the process of drug discovery is of the utmost importance. However, they also acknowledge the challenges associated with translating this knowledge into practical applications. The necessity for a transition from exclusively academic research to the development of phytomedicines is emphasized, along with the establishment of African centers of excellence in drug discovery. This ongoing investigation into the pharmacological properties of these compounds reflects a growing recognition of the therapeutic capabilities embedded in African flora. Lifongo et \emph{al.} \cite{Lifongo2014}, continue this discourse by exploring the bioactivity of Nigerian medicinal plants. The assertion is made that the quality control and standardization of these plant preparations are essential for their acceptance in modern medicine. The authors advocate for a scientific validation of traditional uses to support the development of new drugs, emphasizing the importance of local support for these initiatives. Salmerón-Manzano et \emph{al.} \cite{SalmeronManzano2020}, present a comprehensive overview of global research trends concerning medicinal plants, emphasizing the historical context of their therapeutic properties. The authors underscore the persistent reliance on traditional medicine in underdeveloped regions, where access to synthetic drugs is constrained. However, there is a caveat regarding the unregulated use of these plants, as they have the potential to pose health risks. As posited by Ungogo et \emph{al}. \cite{Ungogo2020}, the challenges of increasing the production of natural compounds are particularly salient in the context of climate change and environmental degradation. The review under consideration emphasizes the necessity of a multidisciplinary approach to drug discovery that fosters collaboration between academia and the pharmaceutical industry, especially in the context of Nigerian bioresources. Adase et \emph{al.} \cite{Adase2022}, provides a comprehensive analysis of Parquetina nigrescens, emphasizing its phytochemical and pharmacological characteristics. Their work emphasizes the necessity of rigorous scientific investigation to ensure the efficacy and safety of herbal medicines, thereby supporting the development of novel drugs for widespread health issues. Finally, Fajinmi et \emph{al.} \cite{Fajinmi2023}, discuss the propagation of medicinal plants in South Africa, highlighting the cultural significance and economic potential of traditional medicine. They advocate for sustainable practices to ensure the availability of these essential resources, which are integral to the healthcare systems in many African communities.
This extensive review of the extant literature showcases the complex and varied role of numerous African medicinal plants within the broad field of healthcare, highlighting their immense potential as valuable sources for new and groundbreaking pharmaceuticals. As illustrated in \Cref{tab:success_plant}, certain aspects indicate that African NPs have the potential to serve as a viable source of pharmaceutical drugs. The text delves into the challenges associated with their use, including critical sustainability issues, the pressing need for rigorous quality control, and the paramount importance of integrating traditional knowledge with modern scientific methods to ensure optimal outcomes. This analysis underscores the necessity for collaborative efforts among researchers, traditional healers, and policymakers to enhance the effectiveness and availability of these extraordinary natural resources.

\begin{center}
	\begin{landscape}
		\begin{table}[htbp]
			\centering
			\caption{Successes and Potential of African Medicinal Plants}
			\label{tab:success_plant}
			\begin{tabularx}{\linewidth}{>{\bfseries}p{5.5cm} X X X}
				\toprule
				\textbf{Aspect} & \textbf{Successes of Natural Products} & \textbf{Untapped Potential} & \textbf{Socioeconomic and Conservation Aspects} \\
				\midrule
				
				\textbf{Drug Discovery and Biological Activity} &
				African medicinal plants have yielded diverse bioactive compounds (e.g., terpenoids, flavonoids, alkaloids, quinones) with antimicrobial, antimalarial, and anti-ulcer properties. &
				Many African plant compounds show favorable drug-like properties and are strong candidates for virtual screening and lead development. &
				Greater focus on preclinical and clinical studies is needed to validate traditional knowledge and accelerate therapeutic discovery. \\
				
				\textbf{Commercialization} &
				Plants like \textit{Harpagophytum procumbens} (devil’s claw), \textit{Pelargonium sidoides}, and \textit{Aloe vera} have found global commercial success. &
				Less than \num{10}\% of over \num{5400} known medicinal species in Africa have been fully explored or commercialized. &
				Sustainable commercialization benefits rural communities and encourages biodiversity through regulated cultivation. \\
				
				\textbf{Antimicrobial and Antimalarial Successes} &
				Compounds from African plants have demonstrated effectiveness against pathogens such as \textit{Helicobacter pylori} and malaria parasites, especially where conventional drug resistance is an issue. &
				Chemical diversity in unexplored species offers new opportunities for antimicrobial innovation. &
				Traditional medicinal knowledge is at risk due to poor documentation; preservation and integration into research is essential. \\
				
				\bottomrule
			\end{tabularx}
		\end{table}
	\end{landscape}
\end{center}

\subsection{Computational methods for natural product analysis}
The integration of computational methodologies has significantly enhanced the study of NPs, facilitating more precise characterization and identification of bioactive substances. Recent advancements in this field include the application of cheminformatics to assess chemical diversity and the utilization of virtual screening techniques to identify potential therapeutic agents from natural sources \cite{MedinaFranco2013}. These advancements facilitate the systematic investigation of extensive natural product databases, ultimately leading to the discovery of new drugs that can more effectively target various diseases. Furthermore, the integration of bioinformatics with conventional approaches has led to a paradigm shift in the realm of drug discovery, enhancing its efficiency and cost-effectiveness. Consequently, researchers can now utilize these tools to identify molecular targets and evaluate the druggability of natural compounds with increased precision. This collaboration between computational techniques and NP research establishes the foundation for innovative therapeutic developments in the future. This methodological evolution not only streamlines the drug discovery process but also increases the likelihood of finding novel compounds that might have been overlooked using traditional methods. By leveraging computational techniques, scientists can more effectively navigate the complexities of natural product chemistry and enhance the development of new therapeutic agents. This transition towards computational methodologies signifies a substantial shift in the manner in which researchers approach natural product drug discovery, thereby facilitating a more informed and strategic approach to identifying promising candidates \cite{Fajinmi2023,MedinaFranco2013}. The future of this field appears promising, with ongoing advancements expected to further refine and improve the effectiveness of natural product-based therapies.

\section{Machine Learning in Drug Discovey}
Machine learning (ML) and deep learning (DL) have rapidly transformed the process of drug discovery. These technologies have enabled the analysis of large, complex datasets and the automation of key steps in the drug development pipeline. These technologies have become indispensable at various stages, from target identification to clinical trials, offering the potential to accelerate timelines, reduce costs, and improve the accuracy of predictions. In the pharmaceutical industry, the processes of drug screening and optimization are of critical importance. The purpose of these procedures is twofold: first, to identify compounds that have therapeutic potential, and second, to ensure the safety and efficacy of these compounds. Recent advancements in ML and DL have significantly enhanced the efficiency and accuracy of these processes, allowing for better candidate selection and optimization strategies. These methodologies employ extensive datasets to predict bioactivity and ADMET properties, thereby facilitating the streamlining of the drug development pipeline \cite{Chen2011,Wang2012,Swanson2024,Guan2019}. The integration of classical machine learning techniques with deep learning applications has been demonstrated to yield more robust and reliable outcomes in the domain of drug discovery.

In this context, significant applications encompass supervised learning methods, which forecast the biological activity and toxicity of compounds, alongside unsupervised learning techniques that identify patterns within extensive datasets. These approaches refine the accuracy of selecting drug candidates and enhance the predictive modeling of drug efficacy and safety outcomes, thereby supporting more informed decision-making throughout the development process. Despite their promise, the use of traditional ML techniques in drug discovery faces several obstacles, such as challenges related to data quality and availability, model interpretability, and computational requirements \cite{YashPatel2019,Ekins2019}. The complex nature of many ML models can lead to a \emph{"black box"} phenomenon, raising concerns about the transparency and reliability of their predictions. Additionally, issues like class imbalance in datasets can hinder the creation of robust predictive models. Overcoming these challenges is crucial to fully harness ML's potential in optimizing drug development and ensuring dependable outcomes in clinical settings. As machine learning continues to advance, it offers both exciting opportunities and significant ethical considerations for the future of drug discovery. The integration of advanced techniques, along with a focus on personalized medicine, holds the promise of transforming therapeutic approaches and improving patient outcomes. It is essential to recognize the importance of ongoing efforts to enhance data integration, transparency, and collaboration within the field. These initiatives are vital for successfully navigating the complex intricacies of this transformative landscape.
 
\subsection{Type of Machine Learning Techniques} 
Traditional ligand-based and docking-based drug discovery approaches have seen their limitations addressed by the advent of ML. ML methods learn from known data points to predict various properties, such as the binding strength between a potential drug candidate and its target protein \cite{Voort2024}. Here is an overview of the type of machine learning used in computational drug discovery (See \Cref{fig:ml_techniques}). 


\begin{figure}[H]
	\centering
	\includegraphics[scale=0.27]{ML_techniques.png}
	\caption{Type of machine learning with drug discovery applications}
	\label{fig:ml_techniques}
\end{figure}

ML techniques employed in drug discovery can be broadly categorized into supervised, unsupervised, and generative learning paradigms, each addressing distinct stages of the drug development pipeline. Supervised learning methods, including classical algorithms such as support vector machines, random forests, and deep neural networks, are widely used for quantitative structure–activity relationship (QSAR) modeling, enabling the prediction of biological activity, toxicity, and pharmacokinetic properties from molecular descriptors and learned representations. Unsupervised learning techniques, such as clustering, dimensionality reduction, and latent space modeling, play a crucial role in chemical space exploration by identifying structural similarities, discovering hidden patterns among compounds, and reducing the complexity of high-dimensional molecular data. Generative learning approaches, including variational autoencoders, generative adversarial networks, and reinforcement learning–based molecular optimization frameworks, facilitate the de novo design of novel antimalarial candidates by learning the underlying distribution of bioactive molecules and proposing chemically valid and synthetically feasible structures.
  
\subsection{Limitations in handling high-dimensional data and non-linearity}

\section{Quantum Computing Fundamentals}
The fundamental principles and concepts of quantum computing, which are distinct from those of classical computing, are predicated on the phenomena of quantum mechanics. Quantum computing utilizes these phenomena to process information in unprecedented ways. The foundation of this technology is rooted in quantum bits, or qubits, which exhibit a fundamental difference from classical bits in that they exist in multiple states simultaneously through a phenomenon known as superposition. This unique property enables quantum computers to perform complex computations with a level of efficiency that far surpasses that of classical systems. Consequently, this technological advancement has led to significant advancements in various fields, including cryptography, optimization problems, and drug discovery. Notwithstanding its potential, the domain of quantum computing confronts substantial challenges, including concerns regarding decoherence, scalability, and error correction, which impede the practical implementation of quantum systems. It is imperative to acknowledge the significance of ongoing research in addressing the prevailing obstacles and realizing the full potential of quantum computing in addressing problems that have historically been considered intractable for classical computers.

\subsection{Quantum computing concepts}

Quantum computing operations generally entail the manipulation of a quantum register through the application of a sequence of pulses that execute unitary transformations based on a designated algorithm. Subsequent to these operations, the quantum state undergoes measurement, resulting in the collapse of the superposition and the production of a single outcome. To ensure the accuracy and reliability of the results, this procedure is typically repeated thousands of times. The outcomes from all measurements are averaged to mitigate quantum uncertainty and achieve statistical precision. It is imperative to reiterate that repetition is of the essence, given that each measurement yields only a probabilistic representation of the quantum system.  The following are some fundamental concepts in quantum computing.

\tcbset{fonttitle=\bfseries\upshape,
	fontupper=\slshape, colback=yellow!32!white,colframe=red!75!black, separator sign dash, breakable, width=1.04\linewidth}
	
\begin{tcolorbox}
	\begin{itemize}
		\item \textbf{\textcolor{codepurple!80}{Quantum bits (qubits) and quantum gates}} : In contradistinction to classical computers, which rely on bits, quantum computers utilize qubits. A qubit is defined as the fundamental unit of information capable of existing simultaneously in a superposition of both 0 and 1. This property of being in multiple states at once is crucial for conducting parallel computations and examining various possibilities concurrently. In contrast to classical bits, which are limited to a binary representation of either 0 or 1, qubits possess the capacity to simultaneously represent both values due to the superposition principle. Quantum gates are operations applied to qubits to manipulate and process quantum information, thereby enabling transformations that alter the probabilities of measuring different qubit states.
		
		\item \textbf{\textcolor{codegreen!85}{Quantum superposition}} : A single qubit, when considered in isolation, offers limited utility. However, it has been demonstrated that this system is capable of performing a significant feat of quantum physic computation by placing the quantum information it holds into a superposition, thereby merging all possible qubit configurations. In the context of quantum physics, when qubits are in a state of superposition, they generate complex, multidimensional computational environments. These environments facilitate the representation of complicated problems in innovative ways. The principle of quantum superposition enables qubits to exist in multiple states simultaneously. This characteristic is employed in quantum algorithms to facilitate large-scale computations by concurrently examining numerous potential solutions.
		
		\item \textbf{\textcolor{codeblue!85}{Quantum entanglement}}: This phenomenon occurs when two or more qubits become interconnected in such a manner that the state of one qubit is dependent on the state of the others. In the context of quantum physics, when qubits are said to be \emph{"entangled"}, this implies that the observation of one immediately influences the state of the other, regardless of their spatial separation. This entanglement facilitates the formation of highly correlated quantum systems and serves as a vital resource in quantum computing algorithms. In the domain of drug discovery, quantum computing employs the principles of superposition and entanglement to conduct intricate simulations, refine molecular structures, and address computational challenges encountered in areas such as molecular dynamics simulations, drug target identification, and drug molecule optimization. By leveraging the capabilities of quantum physics, quantum computing has the potential to transform the drug discovery process and expedite the advancement of personalized medicine.
	\end{itemize}
\end{tcolorbox}
\section{Quantum Machine Learning}

Quantum machine learning (QML) represents a promising frontier at the intersection of quantum computing and artificial intelligence. QML algorithms leverage quantum mechanical phenomena—such as superposition and entanglement—to enhance certain computational tasks within machine learning workflows. The fundamental premise is that quantum systems can process information in ways that classical architectures cannot, potentially offering computational advantages for specific problems in chemistry and drug discovery. Nevertheless, the practical realization of these advantages remains an area of active investigation, with current hardware limitations and algorithm maturity presenting significant challenges.

\subsection{Recent QML algorithms and their potential for accelerating drug discovery}

Several QML algorithmic frameworks have been proposed for applications in computational chemistry and drug discovery. Quantum kernel methods, which leverage the Hilbert space of quantum states to compute similarity measures between molecular representations, have been investigated as alternatives to classical kernel approaches. Variational quantum eigensolvers (VQE) offer a hybrid quantum-classical approach to solving electronic structure problems, potentially enabling more accurate ground-state energy calculations for molecular systems of pharmaceutical interest. Quantum neural networks, composed of parameterized quantum circuits, have been explored for supervised learning tasks such as molecular property prediction.

Despite theoretical promise, empirical evaluations on near-term quantum hardware have yielded mixed results. Recent benchmark studies suggest that for current problem sizes and noise levels, quantum kernel methods do not consistently outperform well-optimized classical baselines such as radial basis function kernels or tree-based ensemble methods. For instance, comparative assessments on molecular activity prediction panels have observed no significant quantum advantage when evaluated under rigorous statistical testing and appropriate multiple comparison corrections. This observation underscores the importance of honest-negative reporting in the field: claims of quantum superiority must be substantiated through controlled experiments with clearly defined metrics, appropriate baselines, and statistical validation.

\subsection{Quantum applications for drug discovery and opportunities for malaria research}

The potential application of quantum computing to malaria drug discovery is predicated on several computational bottlenecks that quantum algorithms may address. Accurate simulation of molecular electronic structure, particularly for transition metal complexes or systems with strong electron correlation, remains computationally expensive for classical methods. Quantum simulation algorithms could, in principle, provide more tractable pathways to these calculations. Similarly, optimization tasks in chemical space exploration—such as identifying molecular candidates that simultaneously satisfy multiple pharmacological constraints—may benefit from quantum annealing or variational optimization approaches.

However, practical deployment of quantum methods for malaria-specific drug discovery faces substantial challenges. Current quantum hardware exhibits limited qubit counts, short coherence times, and high error rates, restricting the size and complexity of molecular systems that can be reliably simulated. Furthermore, the translation of quantum computational results into actionable drug discovery decisions requires integration with experimental validation pipelines, including \emph{in vitro} assays and \emph{in vivo} models, which remain outside the quantum computing paradigm. Consequently, the most promising near-term applications may involve hybrid quantum-classical workflows, where quantum subroutines are invoked for specific subtasks within a predominantly classical computational framework.

\section{Computational Drug Discovery Methods}

The integration of computational techniques into drug discovery has fundamentally transformed the efficiency and scope of therapeutic candidate identification. Computational approaches enable the systematic exploration of vast chemical spaces, the prediction of molecular properties prior to synthesis, and the prioritization of candidates based on multi-objective criteria. These methods complement experimental workflows by reducing the number of compounds requiring physical synthesis and biological testing, thereby accelerating the overall discovery timeline and reducing associated costs.

\subsection{Molecular docking and structure-based design}

Molecular docking is a computational technique that predicts the preferred orientation and binding affinity of a small molecule ligand when bound to a macromolecular target, typically a protein. Docking algorithms employ scoring functions to estimate the free energy of binding, enabling the ranking of candidate compounds based on their predicted affinity for a given target. Structure-based drug design leverages three-dimensional structural information, often derived from X-ray crystallography or cryo-electron microscopy, to guide the identification and optimization of lead compounds.

In antimalarial drug discovery, molecular docking has been extensively applied to key targets within the \emph{Plasmodium falciparum} proteome. These targets include dihydrofolate reductase (PfDHFR), which is essential for folate metabolism; the chloroquine resistance transporter (PfCRT), implicated in drug efflux mechanisms; the ATPase PfATP4, involved in sodium homeostasis; and the protease PfClpP, which plays a role in protein quality control. Computational screening campaigns targeting these proteins have identified numerous candidate compounds with predicted binding affinities in the low micromolar to nanomolar range.

Nevertheless, docking scores must be interpreted with appropriate caution. Scoring functions are approximations of the true binding free energy and exhibit systematic biases depending on the target and chemical series. Docking-predicted affinities do not directly correspond to experimentally measured IC₅₀ or EC₅₀ values, and false positives are common. Validation through orthogonal methods, including surface plasmon resonance, isothermal titration calorimetry, or cellular assays, remains essential. Furthermore, docking studies typically assume a static receptor structure, neglecting conformational flexibility and solvent effects that can significantly influence binding thermodynamics.

\subsection{Molecular dynamics simulations for drug-target interactions}

Molecular dynamics (MD) simulations provide a time-resolved, atomistic description of molecular systems by numerically integrating Newton's equations of motion. In drug discovery, MD simulations are employed to model the dynamic behavior of drug-target complexes, assess binding stability, and estimate free energy differences between bound and unbound states. Force fields, which parameterize atomic interactions, are central to the accuracy of MD simulations. Common choices for biomolecular systems include CHARMM, AMBER, and OPLS, while small molecule parameterization may employ CGenFF, GAFF, or more recent methods based on quantum chemical calculations.

MD simulations can complement docking studies by providing insights into the conformational dynamics of ligand-receptor interactions over nanosecond to microsecond timescales. Free energy perturbation (FEP) and molecular mechanics/Poisson-Boltzmann surface area (MM-PBSA or MM-GBSA) methods enable the estimation of relative or absolute binding free energies, which correlate more closely with experimental affinity measurements than static docking scores. However, MD-based free energy calculations are computationally intensive and require careful protocol design, including adequate sampling, convergence checks, and consideration of protonation states and solvent models.

In the context of antimalarial research, MD simulations have been applied to model drug resistance mechanisms. For instance, mutations in PfDHFR (such as N51I, C59R, S108N, and I164L) and PfCRT (including K76T and K76A) confer resistance to antifolate and chloroquine-based therapies, respectively. MD simulations can elucidate how these mutations alter the binding mode or stability of drug molecules, providing a molecular-level understanding of resistance emergence. Such insights inform the design of next-generation compounds that retain activity against resistant parasite strains.

\subsection{Chemical space exploration and virtual screening}

Chemical space refers to the multidimensional descriptor space in which molecular structures can be represented. Exploration of chemical space aims to identify regions enriched in bioactive compounds while avoiding areas associated with toxicity or poor pharmacokinetics. Virtual screening techniques systematically evaluate large libraries of compounds—often numbering in the millions—to identify candidates with desirable properties. Virtual screening workflows typically integrate molecular fingerprints, such as extended-connectivity fingerprints (ECFP) or molecular access system (MACCS) keys, with similarity-based or machine learning–based models to predict activity.

A critical challenge in virtual screening is achieving chemical diversity while maintaining drug-like properties. Lipinski's rule of five and related heuristics provide guidelines for oral bioavailability, but these rules are not absolute and may exclude biologically relevant chemical scaffolds. Modern approaches employ multi-objective optimization to balance conflicting criteria, such as predicted potency, selectivity, synthetic accessibility, and predicted resistance resilience. Pareto-optimal solutions, which represent trade-offs among these objectives, offer a transparent framework for candidate prioritization.

African natural product libraries represent an underexplored source of chemical diversity for antimalarial drug discovery. These libraries contain scaffolds distinct from synthetic compound collections, often featuring complex polycyclic structures and stereochemical diversity. Computational screening of African medicinal plant–derived compounds has identified numerous candidates with predicted activity against malaria targets. However, translating these computational predictions into validated hits requires addressing challenges related to compound availability, purity, and structural confirmation.

\subsection{Resistance-resilient scoring and polypharmacology}

Drug resistance remains a persistent challenge in malaria treatment, driven by genetic mutations in parasite targets and transporters. Computational methods can integrate resistance-related information into candidate prioritization by evaluating predicted binding affinities not only against wild-type targets but also against clinically relevant mutant forms. A resistance resilience score (RRS) quantifies the ratio of mutant-to-wild-type binding energy, normalized across multiple resistance-associated mutations. Candidates with high RRS values are predicted to retain activity even in the presence of resistance mutations.

Polypharmacology, the concept that a single compound can modulate multiple targets, offers an additional strategy to combat resistance. Multi-target compounds may reduce the selective pressure for resistance emergence and provide synergistic therapeutic effects. Computational polypharmacology workflows evaluate candidate molecules against panels of targets, assessing both on-target activity and potential off-target liabilities. Network-based approaches, such as polypharmacology network scoring (PNS), quantify the connectivity of a compound to multiple targets within a biological network, providing a systems-level perspective on drug action.

Empirical studies have explored correlations between computational polypharmacology metrics and resistance resilience scores. Results indicate that these metrics capture complementary aspects of molecular behavior: RRS reflects target-level robustness to mutations, while polypharmacology metrics assess breadth of target engagement. However, computational predictions of polypharmacology do not constitute evidence of biological efficacy; experimental validation through multi-target assays and phenotypic screens remains necessary.

\section{Molecular Representations for Machine Learning}

The performance of machine learning models in drug discovery critically depends on the choice of molecular representation. Representations encode chemical structures into numerical vectors suitable for input to predictive models. Classical representations, such as molecular fingerprints and physicochemical descriptors, have been widely adopted due to their computational efficiency and interpretability. More recent approaches explore learned representations derived from graph neural networks or sequence-based models, as well as representations informed by quantum chemical or topological considerations.

\subsection{Classical molecular fingerprints and descriptors}

Extended-connectivity fingerprints (ECFP), also known as Morgan fingerprints or circular fingerprints, are among the most widely used molecular representations in cheminformatics. ECFP encodes the local chemical environment of each atom within a specified radius, producing a fixed- or variable-length binary vector. ECFP4 (radius 2) is a common choice, balancing structural resolution with computational efficiency. These fingerprints are particularly effective for substructure searching and similarity-based virtual screening.

Physicochemical descriptors, such as molecular weight, LogP, topological polar surface area (TPSA), and hydrogen bond donor/acceptor counts, provide interpretable features related to drug-like properties. Descriptor-based models are transparent and readily interpretable, facilitating mechanistic insights. However, fixed descriptor sets may fail to capture complex structure-activity relationships, particularly for structurally diverse or novel chemical scaffolds.

\subsection{Graph neural networks and learned representations}

Graph neural networks (GNNs) operate directly on molecular graphs, where atoms are nodes and bonds are edges. Message-passing neural networks (MPNNs), graph isomorphism networks (GINs), and graph attention networks (GATs) iteratively aggregate information from neighboring atoms to produce node-level or graph-level embeddings. These learned representations can capture higher-order structural patterns not explicitly encoded in classical fingerprints.

Despite their architectural sophistication, GNNs do not universally outperform classical fingerprint-based models. Benchmark studies on antimalarial activity prediction panels have demonstrated that ECFP4 combined with random forests achieves competitive or superior performance compared to GNNs, particularly under scaffold-based train-test splits. This result suggests that for moderately sized datasets (on the order of 10⁴ compounds), the inductive bias provided by ECFP's substructure-aware design is highly effective. GNNs may offer advantages in larger-scale datasets or when incorporating auxiliary information, such as protein-ligand interaction data, but these scenarios remain less common in early-stage drug discovery.

\subsection{Topological and quantum-inspired descriptors}

Topological data analysis (TDA) provides a framework for extracting global structural features from molecular graphs through persistent homology. Persistent homology quantifies the multi-scale connectivity and void structures within a molecular graph, producing topological fingerprints (TFP) that encode information orthogonal to classical descriptors. Similarly, quantum-inspired descriptors, such as those derived from quantum kernel methods, embed molecules into Hilbert spaces defined by quantum feature maps.

Empirical evaluations of topological and quantum-inspired descriptors on antimalarial activity prediction tasks have yielded mixed results. Hybrid models combining ECFP4 with topological features achieve modest improvements in predictive performance compared to ECFP4 alone. However, quantum kernel methods do not consistently outperform classical radial basis function (RBF) kernels when evaluated under rigorous cross-validation and statistical testing. These findings underscore the importance of honest-negative reporting: novel representations should be assessed not only for their theoretical novelty but also for their empirical utility relative to well-optimized baselines.

The integration of topological and quantum-inspired features into ensemble models represents a pragmatic approach, leveraging complementary information while maintaining the strong performance of classical representations. However, the marginal gains observed in current studies suggest that these advanced representations may be most beneficial for specific subsets of the chemical space or for tasks requiring fine-grained discrimination between structurally similar compounds.

\section{Generative Models for Molecule Design}

Generative machine learning models offer the capability to propose novel molecular structures that satisfy specified property constraints. These models learn the distribution of known bioactive molecules and generate new candidates by sampling from the learned distribution. Generative approaches are particularly valuable for de novo design, where the goal is to explore uncharted regions of chemical space rather than screen existing libraries.

\subsection{Variational autoencoders and generative adversarial networks}

Variational autoencoders (VAEs) and generative adversarial networks (GANs) are two prominent generative modeling paradigms. VAEs compress molecular representations into a continuous latent space and reconstruct molecules through a decoder network. By sampling from the latent space, novel molecular structures can be generated. GANs employ a generator-discriminator framework, where the generator learns to produce molecules indistinguishable from real training examples.

Scaffold-based VAEs, such as junction tree VAEs or graph-based VAEs, have been developed to ensure that generated molecules adhere to chemical validity constraints. These models decompose molecules into substructures (e.g., rings, functional groups) and reassemble them during generation, reducing the frequency of chemically infeasible outputs. However, generative models face challenges related to mode collapse, where the generator produces a limited diversity of outputs, and the difficulty of incorporating multiple property constraints simultaneously.

\subsection{Reinforcement learning and multi-objective optimization}

Reinforcement learning (RL) frameworks enable the optimization of molecular properties through iterative design cycles. In RL-based molecular design, an agent proposes molecular modifications (e.g., adding or removing atoms, changing bond types) and receives rewards based on the predicted properties of the resulting molecule. Policy gradient methods, such as proximal policy optimization (PPO), guide the agent toward regions of chemical space enriched in desirable compounds.

Monte Carlo tree search (MCTS), combined with RL, provides a systematic approach to exploring the exponentially large space of possible molecular modifications. MCTS balances exploration (trying new modifications) and exploitation (refining promising candidates) through a tree-structured search process. Multi-objective MCTS extends this framework to optimize multiple conflicting objectives—such as predicted potency, synthetic accessibility, and resistance resilience—simultaneously.

Empirical evaluations of MCTS-based molecular design for antimalarial candidates have produced instructive results. Benchmark comparisons against random sampling baselines reveal that, under certain reward landscapes, random exploration can outperform MCTS in terms of scalar reward maximization. This honest-negative finding highlights the importance of reward function design: flat or rugged reward landscapes may not provide sufficient gradient information to guide MCTS effectively. However, Pareto-based multi-objective formulations offer a complementary perspective, providing auditable trade-offs among competing objectives even when scalar optimization does not demonstrate superiority.

\section{Integration of Computational and Experimental Workflows}

The ultimate success of computational drug discovery depends on the seamless integration of computational predictions with experimental validation. Computational models provide prioritization hypotheses, but experimental assays remain the definitive arbiter of biological activity, safety, and efficacy. Effective integration requires bidirectional feedback: computational models inform experimental design, and experimental results refine computational models.

\subsection{Validation cascades and hit-to-lead progression}

A typical validation cascade begins with computational screening to prioritize a manageable subset of candidates for experimental evaluation. Initial experimental assays may include biochemical binding assays (e.g., surface plasmon resonance, fluorescence polarization) to confirm target engagement. Candidates demonstrating acceptable binding affinities progress to cellular assays, assessing activity against parasite cultures. Compounds exhibiting potent cellular activity undergo counter-screens to evaluate selectivity and cytotoxicity against human cell lines.

Hit-to-lead optimization involves iterative cycles of structural modification, computational re-evaluation, and experimental testing. Structure-activity relationship (SAR) studies guide modifications aimed at improving potency, selectivity, and pharmacokinetic properties. Computational methods, including free energy calculations and ADMET predictions, inform which modifications are most likely to yield improved candidates. This iterative process continues until lead compounds meet predefined criteria for advancement to preclinical development.

\subsection{Prospective validation and retrospective analysis}

Prospective validation, where computational predictions are tested in experiments designed after the predictions were made, provides the strongest evidence of model utility. Retrospective validation, where models are evaluated on historical data, is less definitive due to potential overfitting and lack of temporal separation between training and test sets. Rigorous prospective validation studies are resource-intensive but essential for establishing the predictive value of computational workflows.

In the context of antimalarial research, prospective validation campaigns have evaluated computational predictions against experimental outcomes for compound panels not used in model training. These studies assess metrics such as hit rate (the fraction of predicted actives that are experimentally confirmed), enrichment factor (the improvement over random selection), and area under the receiver operating characteristic curve (ROC-AUC). Results vary depending on the target, chemical series, and computational methodology, underscoring the importance of context-specific model evaluation.

\section{Conclusions and Future Perspectives}

This chapter has provided a comprehensive review of the multifaceted landscape of malaria drug discovery, spanning from the biological foundations of the disease to cutting-edge computational methodologies. The persistent global burden of malaria, particularly in sub-Saharan Africa, underscores the urgency of developing novel therapeutic interventions that can overcome emerging resistance mechanisms and provide durable protection for vulnerable populations. The integration of traditional knowledge regarding natural products, particularly African medicinal plants, with modern computational drug discovery approaches offers a promising pathway toward identifying structurally diverse antimalarial candidates.

The emergence of drug resistance represents one of the most formidable challenges in malaria chemotherapy. Mutations in key parasite targets, such as PfDHFR, PfCRT, PfATP4, and PfClpP, have compromised the efficacy of established therapies, necessitating the development of compounds that retain activity against resistant strains. Computational methods, including resistance-resilient scoring frameworks and molecular dynamics simulations of mutant protein-ligand complexes, provide valuable tools for understanding resistance mechanisms at the molecular level and guiding the design of next-generation antimalarials. However, computational predictions must be validated through rigorous experimental testing, including phenotypic assays against resistant parasite lines and pharmacokinetic profiling in relevant animal models.

The exploration of chemical space through computational screening has revealed vast regions of potentially bioactive molecules that remain experimentally untested. Large-scale virtual screening campaigns, leveraging molecular docking, machine learning models, and multi-objective optimization, have identified numerous candidate compounds with predicted activity against malaria targets. Notably, the integration of African natural product libraries into these screening efforts has uncovered scaffolds distinct from conventional synthetic compound collections, offering opportunities for discovering novel modes of action. Nevertheless, the translation of computational hits into validated leads requires addressing challenges related to compound availability, structural confirmation, and synthetic feasibility.

Machine learning and deep learning methodologies have demonstrated considerable utility in predicting molecular properties, prioritizing candidates, and accelerating hit-to-lead optimization. Classical representations, such as extended-connectivity fingerprints, have proven remarkably effective for antimalarial activity prediction, often outperforming more complex learned representations from graph neural networks or transformer models. This observation highlights a critical lesson: model complexity does not guarantee superior performance, and well-optimized classical baselines should serve as benchmarks for evaluating novel methodologies. Topological and quantum-inspired representations offer complementary information and may provide modest improvements when integrated into ensemble models, but their marginal gains must be weighed against increased computational complexity.

Generative machine learning models, including variational autoencoders, generative adversarial networks, and reinforcement learning–based molecular design frameworks, offer the potential to explore uncharted regions of chemical space and propose de novo candidates. Multi-objective optimization approaches, such as Pareto-based Monte Carlo tree search, provide transparent frameworks for balancing conflicting property constraints. However, empirical benchmarks have revealed that in certain reward landscapes, generative approaches may not consistently outperform random sampling baselines. These honest-negative findings underscore the importance of rigorous experimental design, appropriate baseline comparisons, and transparent reporting of both successes and limitations.

Quantum machine learning represents a frontier area with theoretical promise but limited empirical evidence of practical advantage for current drug discovery applications. Quantum kernel methods, variational quantum eigensolvers, and quantum neural networks have been proposed for molecular property prediction and electronic structure calculations. However, benchmark evaluations on antimalarial activity prediction tasks have not demonstrated consistent superiority over classical methods such as radial basis function kernels or tree-based ensembles. Near-term quantum computing hardware faces substantial limitations, including restricted qubit counts, short coherence times, and high error rates, which constrain the size and complexity of molecular systems that can be reliably simulated. Hybrid quantum-classical workflows may offer the most pragmatic pathway for integrating quantum computational methods into drug discovery pipelines.

The successful integration of computational and experimental workflows requires bidirectional feedback, where computational models inform experimental design and experimental results refine computational predictions. Validation cascades, beginning with biochemical assays and progressing through cellular phenotypic screens to in vivo efficacy studies, provide the necessary evidence for advancing candidates toward clinical development. Prospective validation studies, where computational predictions are tested in experiments designed after the predictions were made, offer the strongest evidence of model utility. Retrospective analyses, while informative, are susceptible to overfitting and lack the temporal separation required for definitive validation.

Looking forward, several key opportunities and challenges merit emphasis. First, the systematic integration of resistance-related information into computational screening workflows will enhance the identification of compounds with broad-spectrum activity against wild-type and mutant parasite strains. Second, the expansion of chemical diversity through the exploration of African natural product libraries and the application of generative modeling approaches will enrich the pool of candidates available for optimization. Third, the development of more accurate force fields and free energy calculation methods will improve the predictive power of molecular dynamics simulations, enabling more reliable assessments of binding affinity and resistance resilience. Fourth, the maturation of quantum computing hardware and algorithms may eventually enable quantum methods to address computational bottlenecks that remain intractable for classical approaches, although this prospect remains speculative at present.

Finally, the translation of computational discoveries into clinical therapies requires sustained collaboration among computational chemists, medicinal chemists, biologists, pharmacologists, and clinicians. No single methodology, whether classical or quantum, can independently solve the multifaceted challenges of antimalarial drug discovery. Rather, the integration of complementary approaches, guided by rigorous validation and honest reporting of both positive and negative results, will provide the most robust pathway toward developing the next generation of life-saving therapies. The work presented in this thesis contributes to this integrative vision by applying quantum machine learning and classical computational methods to the identification and optimization of antimalarial candidates, with a focus on rigorous evaluation, transparent reporting, and practical utility for advancing global health outcomes.


